\documentclass[11pt,letterpaper]{article}
% Shared layout. Compile each document from this directory.
\usepackage[margin=0.8in,headheight=14pt,headsep=16pt,footskip=22pt]{geometry}
\usepackage[T1]{fontenc}
\usepackage{lmodern,amsmath,amssymb,booktabs,enumitem,titlesec,fancyhdr,microtype,xcolor,hyperref}
\definecolor{inkblue}{RGB}{25,65,100}
\hypersetup{colorlinks=true,linkcolor=inkblue,urlcolor=inkblue,filecolor=inkblue,pdfauthor={AI-assisted course study notes}}
\setlength{\parindent}{0pt}
\setlength{\parskip}{5pt}
\setlength{\emergencystretch}{1.5em}
\setlength{\abovedisplayskip}{7pt}
\setlength{\belowdisplayskip}{7pt}
\titleformat{\section}{\normalsize\bfseries\color{inkblue}}{}{0pt}{}
\titlespacing*{\section}{0pt}{9pt}{3pt}
\setlist[itemize]{leftmargin=15pt,itemsep=2pt,topsep=3pt,parsep=0pt}
\setlist[enumerate]{leftmargin=17pt,itemsep=3pt,topsep=3pt,parsep=0pt}
\pagestyle{fancy}
\fancyhf{}
\fancyhead[L]{\small ECE 750T\enspace /\enspace Responsible AI}
\fancyhead[R]{\small Fall 2026\enspace /\enspace Study notes}
\fancyfoot[R]{\small\thepage}
\renewcommand{\headrulewidth}{0.3pt}
\newcommand{\lesson}[2]{%
  \hypersetup{pdftitle={ECE 750T: #1 -- #2}}%
  {\LARGE\bfseries #1}\par\vspace{3pt}
  {\large #2}\par\vspace{3pt}
  {\footnotesize AI-assisted reconstruction from course slides, not a lecture transcript.
  Discussion reasoning is inferred; classroom responses were not supplied.}\par
  \vspace{3pt}\hrule\vspace{5pt}}
\newcommand{\sources}[1]{\par\vspace{3pt}{\footnotesize\textbf{Slide trail.} #1
  Page numbers refer to PDF pages, not the slides' continuous printed numbering.}\par}
\newcommand{\Prob}{\mathbb{P}}
\newcommand{\E}{\mathbb{E}}
\newcommand{\Yhat}{\widehat{Y}}
\newcommand{\indep}{\mathrel{\perp\!\!\!\perp}}

\begin{document}
\lesson{Lecture 1}{AI risks, harms, and the system around the model}

\textbf{The lesson.} Responsible AI requires reasoning about the entire decision
system: what is measured, whose goals determine the target, how predictions become
actions, and how those actions change opportunities. The scholarship example makes
these choices visible.

\section*{Start with the decision, then examine the prediction}
A university with limited scholarship funds could predict future grades, financial
need, or previous award decisions. Predicting past awards can reproduce historical
decisions; predicting grades can miss unequal access to resources. Ask what the
scholarship should achieve and what the target leaves out.

The healthcare example gives this concern a concrete mechanism: a system used
predicted healthcare spending to select patients for extra care. Spending also
reflects access to care. At the same score, Black patients had more chronic
illnesses. Accurate cost prediction did not establish comparable measurement of need.

Even an appropriate prediction leaves the action open. The same grade estimate
could rank applicants, be combined with financial need, or trigger academic support.
Assess what happens after the score and who bears each kind of error. Removing a
demographic input or applying one cutoff cannot settle fairness: proxies can encode
group membership, and labels or opportunities may already be unequal.

\section*{Follow the loop, rather than stopping at the test set}
The slides' ML loop runs from the state of the world through measurement, data,
learning, action on individuals, and feedback. Awards change students' resources;
those resources affect later
grades; the grades can become the next training labels. Outcomes among funded
students do not reveal how unfunded students would have performed with support.
The predictor helps create the distribution it later learns from.

Foundation models add a shared dependency: many institutions adapt the same model.
Improvements and inherited blind spots can both spread widely. This
\emph{homogenization} connects local failures to systemic concerns about
discrimination, information quality, privacy, and concentrated decision-making power.
Responsibilities span data providers, developers, deployers, reviewers, and affected people.

\section*{Turn a risk category into a causal account}
For an essay summarizer, a useful harm analysis is:
\emph{reviewers rely on the summary $\to$ relevant context is omitted $\to$ an
applicant loses funding}. Identify the affected person, deployment conditions,
evidence for the failure pathway, and the alternative policy used for comparison.
Listing ``bias'' or ``misinformation'' supplies none of those details.

Risk-based thinking asks what could go wrong and how likely and severe it is.
Harms-based thinking centres who bears costs, accountability, and repair.
Taxonomies help find omissions and connect concerns; a category does not establish
likelihood, severity, or whether a safeguard works.

\newpage
\section*{Agency changes what must be checked}
The scholarship system can \emph{rank} candidates for a committee,
\emph{recommend} awards requiring approval, or \emph{allocate} and release funds
without prior approval. These designs can use the same prediction model while
granting it different authority. The agentic-harms discussion organizes agency
along four dimensions:

\begin{itemize}
\item \textbf{Underspecification:} freedom to choose how to achieve a goal.
\item \textbf{Directness of impact:} ability to affect the world without human intervention.
\item \textbf{Goal-directedness:} strength of pursuit of a quantifiable objective.
\item \textbf{Long-term planning:} linked actions over time toward that objective.
\end{itemize}

These are lenses for inspecting a deployment, rather than a theorem that more
agency always produces more harm. A system's permissions, goal, oversight, and
time horizon determine which failure pathways deserve investigation. An agent
that can act outside its intended environment creates an oversight problem even
when its stated task is benign.

\section*{How to reconstruct the scholarship discussion}
The slides ask whether autonomous allocation should be allowed. A defensible
answer must connect a position to a condition that can actually be assessed.
For example, before allowing allocation, require evidence about errors for relevant
groups, an appropriate target, a way to correct records, and authority to suspend
the system. Each proposed condition should specify whose harm it addresses, what
would demonstrate effectiveness, and who enforces it. These are possible lines
of reasoning, not recorded conclusions from classmates.

``There is a human in the loop'' leaves important questions unanswered: can that
person inspect evidence, overturn an award, detect systematic failures, and act
before funds are released? Likewise, documenting the system or conducting an audit
is a proposed safeguard whose usefulness needs evidence and institutional support.

Individual and collective control are distinct. An appeal challenges one award;
it may offer no route to challenge the objective, budget policy, or provider's
influence. Ask who can change criteria, stop the system, and secure remedy.

\section*{What this establishes for the rest of the course}
Fairness formalizes selected comparisons; evaluation tests behaviour under
specified conditions; measurement asks whether the target represents the intended
construct; privacy constrains information released about contributions to data.
Each supplies evidence for a different part of the deployment decision.

\textbf{Midterm emphasis (inference).} Be ready to analyze a fresh scenario using
\emph{goal $\to$ measurement $\to$ decision $\to$ feedback $\to$ harm}, explain
why prediction does not determine policy, and distinguish a taxonomy from evidence.
The four agency dimensions and the difference between appeal and collective control
are plausible conceptual material. Memorizing course administration, biographies,
or the dates of news examples is a lower priority. The slides explicitly put
Weeks 1--5 in scope; they do not identify which of these concepts will be tested.

\sources{\href{L01.pdf}{L01.pdf}, pp.~24--32 (targets and the ML loop),
33--39 (foundation models and risk lenses), 40--50 (agency and discussion),
52 (synthesis); annotated copy consulted. Course scope: p.~17.
Main readings presented: Barocas, Hardt, and Narayanan; Bommasani et al.;
Chan et al., \emph{Harms from Increasingly Agentic Algorithmic Systems}.}
\end{document}
